\subsection{Planteamiento del Problema}

El Model Context Protocol ha logrado estandarizar la manera en que los LLMs
interactúan con herramientas externas, proporcionando un marco uniforme para el
descubrimiento, invocación y manejo de resultados de herramientas. Sin embargo,
a pesar de sus ventajas, el mecanismo
convencional de \textit{tool calling} sobre MCP presenta limitaciones
estructurales que comprometen su escalabilidad en terminos de costo y eficiencia en aplicaciones de producción
\citep{anthropic2024codeexec}.

El primer problema es la \textbf{carga estática de definiciones de herramientas}
(\textit{static tool loading}). En el paradigma convencional de tool calling,
todas las definiciones de herramientas disponibles incluyendo nombres,
descripciones, parámetros y schemas deben cargarse en la ventana de contexto
del modelo \textit{antes} de que este procese el \textit{prompt} del usuario. Esto
significa que, en escenarios donde un agente está conectado a múltiples
servidores MCP con decenas o cientos de herramientas, el modelo debe procesar
cientos de miles de tokens de definiciones antes de siquiera leer la tarea
por hacer \citet{anthropic2024codeexec}.Se señala que los agentes desperdician contexto significativo en
descripciones que en su mayoría no serán utilizadas para la tarea actual.

El segundo problema es el \textbf{tránsito obligatorio de resultados intermedios
por el contexto del modelo}. En una secuencia de llamadas a herramientas donde
el resultado de una alimenta la entrada de la siguiente, cada resultado
intermedio debe pasar por la ventana de contexto del modelo para ser
interpretado y reformulado como entrada de la siguiente llamada. Esto implica
que los mismos datos son procesados múltiples veces.
\citet{anthropic2024codeexec} ilustran este problema con un ejemplo concreto:
en un flujo de recuperación y actualización de documentos, la transcripción
completa de una llamada fluye por el contexto dos veces; esto significa procesar tokens adicionales
innecesarios. En el caso de documentos extensos o estructuras de datos
complejas, este patrón puede incluso exceder los límites de la ventana de
contexto, haciendo la tarea irrealizable.

El tercer problema es la \textbf{propensión a errores en la manipulación de
datos entre llamadas}. Cuando un modelo debe copiar datos de la respuesta de una
herramienta hacia los parámetros de la siguiente, la probabilidad de errores de
transcripción aumenta proporcionalmente con el volumen y la complejidad de los
datos. \citet{anthropic2024codeexec} advierten que, con documentos extensos o
estructuras de datos complejas, los modelos son más propensos a cometer errores
al copiar datos entre llamadas a herramientas.

Estos tres problemas, carga estática, tránsito redundante de datos y errores
de transcripción comparten una raíz común: el \textbf{la armonía entre la
orquestación de herramientas y la ventana de contexto del modelo}. En el
paradigma convencional, el contexto del modelo cumple una doble función: es
simultáneamente el espacio de razonamiento del LLM y el canal de comunicación
entre herramientas. Esta doble función genera una tensión fundamental entre la
cantidad de herramientas disponibles y la eficiencia con que se utiliza el
contexto.

Anthropic propone como alternativa la ejecución de código,
donde el modelo genera programas que orquestan las llamadas a herramientas en un
entorno externo, descubriendo y cargando definiciones bajo demanda en lugar de
estáticamente. Este enfoque reporta reducciones de hasta 98.7\% en consumo de
tokens de 150,000 a 2,000 tokens en escenarios con grandes conjuntos de
herramientas. Sin embargo, esta cifra proviene de un caso de uso específico y no
constituye una evaluación sistemática a través de diferentes tipos de tareas,
niveles de complejidad y sistemas agenticos.

La pregunta central de esta investigación es, por tanto: \textbf{¿bajo qué
condiciones la ejecución de código a través de MCP supera a las llamadas
directas en términos de consumo de tokens, latencia y precisión en sistemas agenticos?}
